\addcontentsline{toc}{section}{Abstract}
\section*{Abstract}

This thesis presents the design and implementation of \nsit, a web \acrlong{vn} engine built on a declarative \acrshort{json} file. The project stems from a personal observation: existing \acrlong{vn} engines often offer rich features, but impose compilation steps, external dependencies, or deployment pipelines that complicate the task for a non-developer author who simply wants to tell an interactive story. The goal of the project was therefore to test whether a more direct architecture, driven solely by a declarative script interpreted in the browser, could allow this type of author to create and publish interactive content without any installation or compilation on the user's side.

To achieve this, I structured the development around a \glsit{pdca} cycle across three iterations: setting up the engine's foundations, optimising \gls{rendu}, and strengthening robustness. The work resulted in a \glsit{runtime} built in \gls{ts}/\gls{reactjs}, declarative data validation via the \gls{zod} \gls{librairie}, and several mechanisms dedicated to narrative navigation, progressive text display, and the handling of choices and branching.

The results show that it is technically possible to describe an interactive narrative with a simple declarative format and run it directly in a browser, without a \glsit{backend}. The optimisation iteration reduced the number of \glspl{rendu} required by 28.5\%, and validation via \gls{zod} reliably detects \acrshort{json} \gls{syntaxe} errors as well as corrupted save files. However, the \acrshort{json} format remains demanding for a non-technical author, which confirms that an engine of this kind should ideally be complemented by a guided editing tool to be truly accessible to non-developer authors.

\medskip
\noindent\textbf{Keywords:} \textit{Visual Novel}, declarative script, JSON, no-code/low-code, interactive narration.
